\section{Exact roof laws at all interval resolutions}
\label{sec:v56-resolution-laws}
We apply the same positive maximal envelope to the path-valued raw
measure. It gives estimates for every subinterval containing a good
roof, without prescribing how fast the subinterval shrinks. The
conditioning event remains the original exact-label event with that
subinterval. A statement uniform over intervals is not an assertion
that conditioning on the output of an adaptive interval-selection
procedure is the same event.

Let $\mathbf P(u)$, $\mathbf b_B(u)$ and $\mathsf W_R$ be the
positive collision-bridge kernels and Gaussian reference of
Section~\ref{sec:v52-path-inversion}. Their masses are $P(u)$ and
$V_B(u)$, respectively. All versions are fixed by its common source
disintegration. Put
\begin{equation}\label{eq:v56-path-discrepancy}
 \eta^{\rm path}_{B,m}
  =\sup_{R,n,k}\operatorname*{ess\,sup}_u
       \|\mathbf P(u)-G(u)\mathsf W_R-\mathbf b_B(u)\|_{\mathrm{BL}^*}.
\end{equation}
Theorem~\ref{thm:v52-path-remainder} gives
$\limsup_m\eta^{\rm path}_{B,m}\le CB^{-1/192}$ for each fixed
$B$. The constant test one belongs to the norming class, so the
scalar discrepancy $\eta_{B,m}$ is at most
$\eta^{\rm path}_{B,m}$. This is a path bounded-Lipschitz dual norm,
not variation mass on path space.

\subsection{One physical envelope for a continuum of resolutions}
\begin{theorem}[Path and scalar control on all anchored intervals]
\label{thm:v56-all-resolutions}
Fix $h>0$ and a sufficiently large band $B$. Use exactly the open
set $O_{B,m,R}^{n,k,I}$ of \eqref{eq:v56-band-threshold}, and put
$r_{B,m}=\lambda_B+\eta^{\rm path}_{B,m}$.
For every $u\in I\setminus O_{B,m,R}^{n,k,I}$,
\begin{equation}\label{eq:v56-all-resolution-path}
 \sup_{J\subset I:\,u\in J}\frac1{|J|}\int_J
       \|\mathbf P(v)-G(v)\mathsf W_R\|_{\mathrm{BL}^*}\,dv
                                                    \le r_{B,m}.
\end{equation}
Almost everywhere on this same complement the integrand at $u$
is at most $r_{B,m}$. The supremum includes every positive interval
length; the open exceptional set does not depend on $J$ or on a
path test. The estimates \eqref{eq:v56-ordered-length} and
\eqref{eq:v56-ordered-source} hold for this set, and
\[
 \limsup_{m\to\infty}r_{B,m}\le CB^{-1/384}.
\]
The same assertions hold for the scalar absolute error and for
integrals against measurable roof-dependent unit bounded-Lipschitz
path tests.
\end{theorem}
\begin{proof}
Positivity gives $\|\mathbf b_B(v)\|_{\mathrm{BL}^*}=V_B(v)$.
Hence, on the common version,
\[
 \|\mathbf P(v)-G(v)\mathsf W_R\|_{\mathrm{BL}^*}
                    \le V_B(v)+\eta^{\rm path}_{B,m}
\]
for almost every $v$. Average over an interval containing a good
anchor and use \eqref{eq:v56-every-interval}. The pointwise assertion
follows from $V_B\le\mathcal M_I V_B$ almost everywhere, not from
an interchange of a collision limit and differentiation. The source
bounds were already proved for this very set. A measurable test
whose restriction to each roof is in the unit path class is bounded
by the displayed dual norm under the integral. One common countable
norming class supplies the roof versions; no uncountable union of
exceptional null sets is taken.
\end{proof}

For a subinterval $J$ of positive length let
\begin{equation}\label{eq:v56-interval-posteriors}
 F_J=\int_JP,\quad H_J=\int_JG_+,\qquad
 d\pi_J=P\one_J\,du/F_J,\quad
 dQ_J=G_+\one_J\,du/H_J.
\end{equation}
These are conditional roof laws only when their denominators are
positive. We use probability total variation
$d_{\rm TV}(\pi,Q)=\tfrac12\int|d\pi-dQ|$.
No factor $1/2$ is inserted in a variation mass of an unnormalized
source.

\begin{corollary}[Conditional laws on arbitrarily short anchored intervals]
\label{cor:v56-interval-posteriors}
Let $u\in I\setminus O_{B,m,R}^{n,k,I}$, let $J\subset I$
contain $u$, and assume $H_J\ge d|J|$ for a fixed $d>0$.
If $r=r_{B,m}<d$, then
\begin{align}
 F_J&\ge(d-r)|J|,\qquad
 d_{\rm TV}(\pi_J,Q_J)\le\frac{r}{d-r},
                                      \label{eq:v56-window-TV}\\
 \int_Jd_{\rm BL}(\mathsf Q^{n,k,v}_{m,R},\mathsf W_R)
                   \,\pi_J(dv)&\le\frac{2r}{d-r}.
                                      \label{eq:v56-window-path}
\end{align}
The joint roof--path law is also within $2r/(d-r)$ of
$Q_J\otimes\mathsf W_R$ when tested by functions measurable in
the roof and unit bounded-Lipschitz in the path. These estimates
hold simultaneously for all such $J$, including intervals depending
on $m$ with no positive lower resolution.

If $G(u)\ge d$ at a good roof, its original same-roof collision
bridge satisfies, almost everywhere there and for
$\eta^{\rm path}_{B,m}<d$,
\begin{equation}\label{eq:v56-good-roof-bridge}
 d_{\rm BL}(\mathsf Q^{n,k,u}_{m,R},\mathsf W_R)
                 \le\frac{2r_{B,m}}{d-\eta^{\rm path}_{B,m}}.
\end{equation}
If $G\ge d$ almost everywhere on $J$ and $r\le d/2$, then the
likelihood $r_J=d\pi_J/dQ_J$ also obeys
\begin{equation}\label{eq:v56-good-roof-likelihood}
 |r_J(u)-1|\le4r/d
\end{equation}
almost everywhere at good roofs $u\in J$. This is not an
essential-supremum estimate on the whole of $J$.
\end{corollary}
\begin{proof}
Write $\mathbf P(v)=P(v)\mathsf Q^{n,k,v}_{m,R}$. For $G(v)<0$,
positivity implies
$\|\mathbf P(v)-G(v)\mathsf W_R\|_{\mathrm{BL}^*}=P(v)+|G(v)|$.
Thus replacing $G$ by $G_+$ never increases this norm. In particular,
\[
 \int_J|P-G_+|\le r|J|,\qquad
 \int_J\|\mathbf P-G_+\mathsf W_R\|_{\mathrm{BL}^*}\le r|J|.
\]
It follows that $|F_J-H_J|\le r|J|$ and $F_J\ge(d-r)|J|$.
Adding and subtracting $G_+/F_J$ bounds the roof variation mass
by $2r|J|/F_J$. Dividing by two proves
\eqref{eq:v56-window-TV}. Normalizing the path-valued inequality
costs at most a further $|F_J-H_J|/F_J$, proving the joint-test
assertion. Also, pointwise,
\[
 P(v)d_{\rm BL}(\mathsf Q^{n,k,v}_{m,R},\mathsf W_R)
       \le2\|\mathbf P(v)-G(v)\mathsf W_R\|_{\mathrm{BL}^*}.
\]
Integration proves \eqref{eq:v56-window-path}.

At a good roof the unnormalized dual error is at most $r$.
The scalar positive-remainder identity gives
$P(u)\ge G(u)-\eta^{\rm path}_{B,m}$; division proves
\eqref{eq:v56-good-roof-bridge}. For the last assertion $G_+=G$
on $J$, and the exact identity is
\[
 r_J(u)-1=\frac{H_J(P(u)-G(u))}{F_JG(u)}
                              +\frac{H_J-F_J}{F_J}.
\]
Here $H_J/F_J\le1+r/(d-r)\le2$ and $|P(u)-G(u)|\le r$ at
almost every good roof. The two terms are each at most $2r/d$.
The pointwise floor is used exactly here; an integrated lower bound
alone would not justify the likelihood claim.
\end{proof}

\subsection{The original roof law sees few bad anchors}
\begin{corollary}[Probability of the common exceptional set]
\label{cor:v56-anchor-probability}
On a mother interval $I$ with $\int_I G_+\ge dh$, the conditional
roof laws of \eqref{eq:v56-interval-posteriors} satisfy
\begin{equation}\label{eq:v56-bad-probability}
 \limsup_{m\to\infty}\sup_{R,n,k,t}
 \{\pi_I(O_{B,m,R}^{n,k,I})+Q_I(O_{B,m,R}^{n,k,I})\}
                                 \le C_{h,d}B^{-1/384}.
\end{equation}
There are sequences $B_m\to\infty$ and $a_m\to0$, independent of
the exact labels and translated mother interval, such that the
errors in Theorem~\ref{thm:v56-all-resolutions} and
Corollary~\ref{cor:v56-interval-posteriors} tend uniformly to zero
outside sets of Lebesgue, $\pi_I$ and $Q_I$ measure tending to zero.
The errors are simultaneous over all anchored subintervals. No
polynomial rate in $m$ is asserted for $B_m$ or $a_m$.
\end{corollary}
\begin{proof}
The inherited local $L^1$ law gives $\int_I P\ge dh/2$ eventually.
Divide \eqref{eq:v56-ordered-source} by this bound. For the reference
use its density $G_+/\int_I G_+\le C_G/(dh)$ and
\eqref{eq:v56-ordered-length}. This proves
\eqref{eq:v56-bad-probability} without assuming that bad anchors
are independent of the source.

For clarity, the diagonal choice uses only uniform fixed-band limits.
Choose $B_j\uparrow\infty$ such that $CB_j^{-1/384}\le2^{-j}$.
For each $j$, choose increasing integers $M_j$ so that for every
$m\ge M_j$ the path discrepancy at this fixed $B_j$ and the mother
normalization error satisfy their corresponding upper bounds with
slack $2^{-j}$. Set $B_m=B_j$ for $M_j\le m<M_{j+1}$.
The finite length bound and the finite source bound in the proof of
Corollary~\ref{cor:v56-ordered-localization} then give the stated
vanishing exceptional measures. The all-interval inequality is
already finite for each $(B,m)$, so it requires no further diagonal
over subinterval lengths. For example $a_m=C_{h,d}2^{-j}$ on these
blocks bounds the errors once $j$ is large. Constants at growing
spectral bands have not been estimated or substituted into a rate.
\end{proof}

\subsection{All-resolution transfer for the actual return bridge}
The physical maximal set above is quantitative in the ordered band
limit for the collision bridge. The actual-return clock transfer is
available in an integrated norm. The following separate application
makes its quantifiers explicit and does not assign it the preceding
band exponent.

Let
\[
 E_m^{\rm ret}(v)=
 \|P(v)\mathsf R^{n,k,v}_{m,R}-G(v)\mathsf V_R\|_{\mathrm{BL}^*},
 \qquad E_m^{\rm col}(v)=
 \|\mathbf P(v)-G(v)\mathsf W_R\|_{\mathrm{BL}^*}.
\]
The target class is $|Z(t)|\le M_0$ with fixed $h,M_0$, using the
prescribed enlargement $|Z(v)|\le M_0+h/\sqrt m$ on $I$.

\begin{theorem}[A common all-resolution set for both actual bridges]
\label{thm:v56-return-maximal}
There is a deterministic $e_m\downarrow0$ such that, uniformly in
this central class,
$\int_I(E_m^{\rm ret}+E_m^{\rm col})\le(1+h)e_m$.
Set $\ell_m=\sqrt{e_m}+m^{-1}$ and
\[
 O_m^{\rm both}=
 \{\mathcal M_I(E_m^{\rm ret}+E_m^{\rm col})>\ell_m\}.
\]
Then
\begin{align}
 |O_m^{\rm both}|&\le C(1+h)e_m/\ell_m,
                                      \label{eq:v56-return-bad-length}\\
 \int_{I\cap O_m^{\rm both}}P
       &\le C(1+h)(e_m/\ell_m)^{1/145}.
                                      \label{eq:v56-return-bad-source}
\end{align}
For every anchor outside this set and every subinterval $J\subset I$
containing it, both integrated dual errors divided by $|J|$ are at
most $\ell_m$. The following conditional bounds are for $\ell_m<d$,
which holds eventually for each fixed $d>0$.
If $\int_JG_+\ge d|J|$, both original conditional
bridges obey the joint-test and conditional-mean estimates of
Corollary~\ref{cor:v56-interval-posteriors}, with $r=\ell_m$.
At almost every good roof where $G\ge d$, their same-roof
bounded-Lipschitz distances are at most $2\ell_m/(d-\ell_m)$.
These conclusions are uniform over every positive subinterval length,
not over arbitrary externally chosen roofs inside the exceptional set.
\end{theorem}
\begin{proof}
The first assertion follows from
Lemma~\ref{lem:v52-return-transfer} and
Theorem~\ref{thm:v52-path-remainder}; replace the suprema of their
mean errors by decreasing tail suprema. Their errors are bounded
by $2(P+|G|)$ in sum, so their local integrability and common roof
versions cause no difficulty. Apply
\eqref{eq:v56-maximal-weak-one} to that sum and then
\eqref{eq:v56-source-weight} to $P$. This proves both set bounds.
All interval averages at a good anchor are at most $\ell_m$ by
definition, and both pointwise errors have that bound almost
everywhere there by differentiation. The constant path test yields
$|P-G|\le\ell_m$ at such roofs. Repeat the positive-part and
normalization argument of Corollary~\ref{cor:v56-interval-posteriors}
for each bridge separately. Since $P\ge G-\ell_m$, the claimed
same-roof denominator follows. No return event, clock value or source
law is replaced, and no quantitative rate is assigned to $e_m$.
\end{proof}

\paragraph{Relation to the full pointwise endpoint.}
The distinction from Corollary~\ref{cor:v52-full-roof-sets} is the
supremum over all intervals containing the same good anchor. That
corollary controls the kernel at most roofs but places no simultaneous
bound on every arbitrarily short surrounding interval. The covering
argument supplies this uniformity directly from the unchanged positive
source or the already proved integrated return numerator. It is a
real-variable consequence of those dynamical inputs, not a new
spectral theorem or a claim of a second singular-hyperbolic application.
The local limit results of \cite{SV,DPZ,DN} supply the established
background; none is differentiated here into an unproved density bound.

The common sets are target- and count-dependent. All densities and
conditional laws in the conclusions still include the full physical
source, including its values on these sets. Ruling out an upward spike
inside them remains precisely the positive incidence and clearance
height problem of Corollary~\ref{cor:v51-positive-criterion}.
A vanishing bound on their measure is not a proof that they are empty.
